\documentclass[12pt, a4paper]{article}

% --- Pacotes Básicos ---
\usepackage[utf8]{inputenc}
\usepackage[]{graphicx}
\usepackage[T1]{fontenc}
\usepackage{times}
\usepackage{setspace}
\onehalfspacing % Set line spacing to 1.5
\setlength{\parindent}{1.25cm}
\setlength{\lineskip}{1.5cm}
%\usepackage[main=brazilian, english]{babel}
\usepackage{indentfirst}
\usepackage{geometry}
\geometry{left=3cm, top=3cm, right=2cm, bottom=2cm}

\usepackage{etoolbox} % Pacote necessário para ajustar o espaçamento

% Criando o ambiente 'citacao'
\newenvironment{citacao}{
  \vspace{0.5cm} % Espaço antes da citação
  \small % Letra menor (geralmente tamanho 10)
  \begin{list}{}{
    \leftmargin=4cm % Recuo de 4 cm da margem esquerda
    \rightmargin=0cm
    \parsep=0pt
    \topsep=0pt
    \partopsep=0pt
  }
  \item\relax
  \linespread{1.0}\selectfont % Garante espaçamento simples dentro da citação
}{
  \end{list}
  \vspace{0.5cm} % Espaço depois da citação
}
\usepackage{fancyhdr}

% Configura o estilo "fancy"
\fancypagestyle{plain}{
  \fancyhf{} % Limpa todos os cabeçalhos e rodapés
  \fancyfoot[R]{\thepage} % Define a paginação no rodapé à direita (R = Right)
  \renewcommand{\headrulewidth}{0pt}  % Remove a linha divisória do cabeçalho
  \renewcommand{\footrulewidth}{0pt}  % Remove a linha divisória do rodapé
}

% Ativa o estilo configurado para todo o documento
\pagestyle{plain}
% --- Informações do Documento ---
\title{UNIVERSIDADE FEDERAL DA FRONTEIRA SUL\\Ciência da Computação\\Meio Ambiente, Economia e Sociedade\\Prof. Leandro Bordin\\ Movimento Ambientalista: Desenvolvimento Sustentável em Meio ao Capitalismo}
\date{}

\begin{document}
\maketitle

\begin{flushright}
	Erickson Giesel Muller
\end{flushright}

\section{A Irreversibilidade e as Dimensões da Equidade}

	O método de produção do capitalismo exploratório no século XXI vem causando grandes impactos ambientais e sociais. Quando falamos em sustentabilidade, pensamos em uma abordagem ecológica com abrangência de equidade intrageracional (entre as gerações atuais), intergeracional (com as gerações futuras) e internacional (com a população de outros países) (Eden, 1994). A visão que o método de produção atual tem reflexos nessas três esferas é de suma importância para a garantia do direito ao meio ambiente, contudo, dentro de um sistema que exige expansão contínua de mercado, deve-se atentar aos colaterais que vão além dos reflexos ecológicos.
\\

	O desenvolvimento econômico é marcado pelo conflito entre melhorias sociais em contrapartida da irreversibilidade de seus impactos ambientais. Segundo Sandroni, o desenvolvimento sustentável é "conceito que pertence ao campo da ecologia e da administração e que se refere ao desenvolvimento de uma empresa, ramo industrial, região ou país, e que em seu processo não esgota os recursos naturais que consome nem danifica o meio ambiente de forma a comprometer o desenvolvimento dessa atividade no futuro." Contudo, diante de uma máquina empresarial que visa lucros acima de tudo, a adoção de uma produção minimamente ecológica virou peça de \textit{marketing} para essas indústrias, que acabam dispensando a preocupação com os impactos sociais.

\begin{citacao}
	DESENVOLVIMENTO ECONÔMICO. Crescimento econômico (aumento do Produto Nacional Bruto per capita) acompanhado pela melhoria do padrão de vida da população e por alterações fundamentais na estrutura de sua economia. O estudo do desenvolvimento econômico e social partiu da constatação da profunda desigualdade, de um lado, entre os países que se industrializaram e atingiram elevados níveis de bem-estar material, compartilhados por amplas camadas da população, e, de outro, aqueles que não se industrializaram e por isso permaneceram em situação de pobreza e com acentuados desníveis sociais. (SANDRONI, 1999, pg. 169)
\end{citacao}

	O problema não é apenas "salvar a natureza", mas lidar com o esgotamento físico e a degradação social. Nesse cenário entram os movimentos ambientalistas, que buscam promover o ecodesenvolvimento com o objetivo de aproximar a igualdade social. Segundo Diani, "um movimento social não é um mero conjunto de ativistas, mas  redes de interação informal entre uma pluralidade de indivíduos, grupos, e organizações, engajados em conflitos políticos ou culturais, com base em identidades coletivas compartilhadas". O papel dos movimentos ambientalistas é pressionar o mercado pela adoção de tecnologias sociais para absorver custos sociais, não apenas os ambientais.

\section{As Três Lentes Econômicas}
% Três correntes da economia ambiental: economia ambiental neoclássica, economia ambiental ecológica e economia ambiental marxista.
	Existem três correntes da economia ambiental, cada uma enxerga a solução para a crise socioambiental de uma maneira distinta: economia ambiental neoclássica, economia ambiental ecológica e economia ambiental marxista. A corrente neoclássica é a visão dominante no mercado atual, enxerga a degradação ambiental como uma falha de mercado, ou seja, uma externalidade negativa. Nessa perspectiva, a palavra "recursos" é primeiro nome identificador do termo "recursos naturais", segundo Montibeller, a economia ambiental neoclássica "parte do pressuposto de que toda externalidade, isto é, todo recurso ou serviço ambiental não incluído no mercado, pode receber uma valoração monetária convincente".

	\begin{citacao}
		A postura da teoria neoclássica, como de quase toda a ciência, desconsiderando ou não tomando como elemento ativo a componente ambiental, até por volta dos anos 1960, explica-se pelo fato de a pressão das atividades econômicas e das concentrações geográficas sobre o meio ambiente não terem atingido, até este marco histórico, grau capaz de gerar efetiva consciência ecológica na sociedade. (Montibeller, 2004, pg. 58)
	\end{citacao}

	Essa abordagem define um valor econômico para os recursos naturais, para a visão da indústria esse valor se faz como custo privado e presente, enquanto, para a sociedade, o custo desses recursos naturais é consideravelmente maior, tendo em vista os reflexos irreversíveis e de impacto social da degradação ecológica. Enquanto para as empresas esses recursos são benefícios quase que gratuitos, a sociedade segue custeando esse benefício sem um retorno à altura.
\\
	
	As empresas tomam decisões baseadas apenas nos custos privados, descartando a análise dos custos ambientais. Nesse sentido, os movimentos ambientalistas se fazem presentes para pressionar a empresa para que esta assuma, não apenas os custos privados, mas também os custos sociais que antes não eram assumidos (logo, faziam-se custos socializados). Os movimentos ambientalistas tomam partes como representantes de interesses nessa troca de valor econômico.
	\begin{citacao}
		Diz-se acima ser factível se chegar a um valor econômico que incorpore todas as formas de valor - isto é, o valor de uso, o valor de opção e o valor de existência - de um bem ambiental; mas que, porém, a avaliação não é correta. A representação dos não-humanos e das gerações futuras é feita pelo avaliador atual e depende, portanto, do caráter altruísta, não egoísta, do indivíduo que manifesta as preferências. Então, por princípio, a representação dos interesses que não diretamente os seus é precária, e os pesos que atribui aos interesses alheios é inferior ao que deveria ser. (Montibeller, 2004, pg. 89)
	\end{citacao}

	Diferente da visão neoclássica, a economia ecológica é interdisciplinar e trabalha com as leis da física, deixa de ser um sistema isolado e passa a ser um subsistema da biosfera. A irreversibilidade dos danos ecológicos levada em conta na tomada de decisão das indústrias, que buscam um equilíbrio entre desenvolvimento econômico e manutenção do ecossistema que abrange populações humanas e não-humanas.
	\\

	Os recursos do planeta são finitos, portanto não se fala mais em "precificar" a natureza, mas em limitar o uso de recursos. O ambientalismo defende a redução do consumo, a transição energética rigorosa e a adoção de métricas que vão além do PIB para medir o desenvolvimento. Neste ambiente, dois conceitos são destacados: a ecologia humana e a economia ecológica. A ecologia humana estuda o uso de recursos em ecossistemas habitados pelo ser humano, definida como ciência de caráter interdisciplinar que analisa os processos biológicos, físicos e sociais que se dão entre os homens e o ambiente em que vivem, segundo Wolanski (1990). A economia ecológica usa os estudos da ecologia humana para compor a análise do desenvolvimento econômico.
	\begin{citacao}
		A economia ecológica, ou ecoeconomia, analisa a estrutura e o processo econômico de geossistemas sob a ótica dos fluxos físicos de energia e de materiais. Trata de explicar o uso de materiais e energia em ecossistemas humanos, mas vai além da ecologia humana, pelo fato de integrar na análise desses fluxos a crítica aos mecanismos e preços de mercado e à valoração econômica da economia ambiental neoclássica. (Montibeller, 2004, pg. 106)
	\end{citacao}

	As críticas dos ecoeconomistas à economia ambiental neoclássica referem-se ao fato desta considerar a variação da base de preços no mercado antes de considerar o fluxo de recursos naturais. Também quanto ao aspecto da "incomensurabilidade de valores", que, segundo Montibeller, significa a "inexistência de padrão de medida comum com outra grandeza", o mercado não é capaz de prever o futuro para poder precificar no presente os recursos cuja falta serão de grande dano às gerações futuras. Nesse sentido, a prioridade pela equidade intergeracional dos ecoeconomistas excede os neoclássicos.
	\\

	A economia ambiental marxista, ou ecomarxismo, aponta o grande conflito entre as vontades de preservação ambiental e manutenção do capitalismo. Baseada nas relações entre valor e mais-valia, o capital humano se enfraquece diante da vontade de geração de margem de lucro pelas indústrias. Para além do funcionamento interno da economia, a produção capitalista acarreta em custos externos ou custos sociais. Segundo Montibeller, "Custos sociais são, portanto, recursos não mercatáveis usados pelo mercado". Os custos sociais são inerentes ao processo produtivo capitalista, a natureza está inclusa como matéria-prima na fabricação de qualquer mercadoria.Os custos externos, as não-mercadorias que compõem as mercadorias, são socializados para fora das portas físicas e temporais da empresa que fabricou essa mercadoria.
	\\

	Com a globalização, o capitalismo selvagem permite que nações superdesenvolvidas exportem seus colaterais ambientas para o Sul Global, transfere o ônus da degradação para as economias periféricas que arcam com o esgotamento dos recursos naturais e a contaminação  de seus territórios para alimentar os padrões de consumo do Primeiro Mundo. Enquanto o Norte Global apresenta uma vitrine de produção de tecnologias "verdes", a periferia é subordinada à condição de fornecedora de matéria-prima barata e destino final de resíduos sólidos, consolidando uma disparidade onde a sustentabilidade do Primeiro Mundo é sustentada pelo subdesenvolvimento socioambiental do Terceiro Mundo.
	\\

	O movimento ambientalista assume um caráter de luta de classes e anti-imperialista. A pressão não é apenas para melhorar o mercado, mas para transformar as estruturas de poder. Nesse sentido, a pauta ambiental transcende a gestão técnica de recursos e passa a ser uma reivindicação pela soberania dos povos sobre seus territórios. O movimento ambientalista demonstra que a  crise ecológica não é uma falha técnica, mas uma decorrência planejada da força de expansão infinita da indústria sobre um planeta finito.

\section{O Papel do Movimento Ambientalista}
	
	Pode-se perceber que o papel do movimento ambientalista contemporâneo vai além da proteção e preservação da natureza, mas se constitui como ator político no tensionamento das relações de produção. Desvelando a noção do crescimento econômico infinito, os movimentos ambientalistas exigem que o capital abandone a relação com a natureza de maneira a tratar desta como fonte de recurso inesgotável e gratuito. Ao fazer isso, o ativismo ambiental expõe as fraturas do sistema, demonstrando que o lucro corporativo frequentemente se sustenta na socialização dos danos ao meio ambiente.
	\\
	
	A internalização das externalidades é fator pressionador em defesa da sustentabilidade no capitalismo ou além deste. O mundo, além de precisar ser um ambiente habitável, precisa que seus habitantes estejam em condições saudáveis para poder ocupar este espaço. Por conta disso, as equidades intergeracionais, intrageracionais e internacionais são ponto focal nas reivindicações dos movimentos ambientalistas. A meta é forçar o mercado a incorporar o custo real da exaustão dos recursos naturais no preço final das mercadorias, reduzindo os impactos destrutivos da atividade econômica e disseminando um padrão de produção que tem o ser humano e a ecologia como nodo central.






\begin{thebibliography}{99}
	\bibitem{SANDRONI} SANDRONI, Paulo. Novíssimo dicionário de economia. São Paulo: Best Seller, 1999.
	\bibitem{MONTIBELLER} MONTIBELLER FILHO, Gilberto. O mito do desenvolvimento sustentável. 2. ed. Florianópolis: Editora da UFSC, 2004.
	\bibitem{EDEN} EDEN, Sally E. Using sustainable development: The business case. Global Environmental Change, v. 4, n. 2, p. 160-167, 1994.
	\bibitem{DIANI} DIANI, Mario. Green networks. A structural analysis of the Italian environmental movement. 1995.
	\bibitem{WOLANSKI} WOLANSKI, Napoleon. A Ecologia Humana. Palestra do autor em simpósio internacional de ecologia humana, Madrid, 1986. Artigo traduzido por Stefan Bulawski Departamento de Filosofia UFSM, 1990.
\end{thebibliography}
%
\end{document}
